\section{Is the Gap Knowledge? An 8B LLM Classifier}\label{sec:llm}

The encoder track and the error analysis (\S\ref{sec:analysis}) show that the model, and not label noise, is what fails.
That leaves two explanations.
The encoder may lack \emph{knowledge} that a 0.4B-parameter model never acquired, or it may lack \emph{label meaning}, which a classifier has to infer from about 3,400 examples.
A larger pretrained classifier tests the first explanation with nothing else changed.
Before any run we registered three predictions.
The gain would pass the screening bar, the single-model \Stwo{} mean would lie in 0.43--0.50 with \Sone{} rising by +0.02 to +0.05, and the gain would fall mostly on the commitment boundary (Explicit $\leftrightarrow$ Implicit, General $\leftrightarrow$ Implicit/Explicit) and more on agreed than on contested items.

\paragraph{One change: the backbone.}
We took the best encoder configuration (sub-question, full question and answer, 16 epochs) and replaced DeBERTa-v3-large with Qwen3-8B-Base \citep{qwen3}.
The model is adapted with LoRA \citep{hu2022lora} of rank 16 ($\alpha=32$, dropout 0.05) on all linear projections of the attention and feed-forward blocks, and a new 9-way linear head, trained in full, reads the hidden state of the last token (Figure~\ref{fig:arch}).
The input text and token budgets are the encoder's, with one fixed closing question appended (``How does the answer respond to the sub-question?'') so that the token the head reads comes after the whole question and answer in every item.
Every run first checks that an item's logits are the same alone and inside a right-padded batch, which guards against a head that reads a padding token.
Training used AdamW at $10^{-4}$ for 3 epochs with an effective batch of 16, and the learning rate and epoch count change only because the backbone does.
Everything else stays identical, namely the same 3,103 training rows, the same 345-row train slice (asserted identical in code), plain cross-entropy and the epoch with the best slice macro-F1.
Seeds 0--9 are paired with the DeBERTa runs of the same seeds, and Appendix~\ref{sec:app} lists all hyperparameters and the compute.

\begin{table*}[t]
\centering
\small
\setlength{\tabcolsep}{5pt}
\begin{tabular}{@{}lccccc@{}}
\toprule
Dev macro-F1 & Qwen3-8B + LoRA & DeBERTa-v3-large & $\Delta$ & Seeds up & 95\% CI \\
\midrule
Single model, \Stwo{}            & \pmsd{0.476}{0.043} & \pmsd{0.384}{0.030} & +0.092 & 10/10 & \\
Single model, \Sone{}            & \pmsd{0.709}{0.031} & \pmsd{0.614}{0.027} & +0.095 & 10/10 & \\
Single model + LA, \Stwo{}       & \pmsd{0.538}{0.043} & \pmsd{0.389}{0.019} & +0.149 & 10/10 & \\
System, \Stwo{}                  & 0.543 (0.549) & 0.405 (0.412) & +0.136 & & [+0.054, +0.224] \\
System, \Sone{}                  & 0.746 & 0.648 & +0.098 & & \\
System, \Stwo{}, 2-annotator sets & 0.524 & 0.382 & +0.142 & & \\
\midrule
Follow-up (one change)           & Variant & Qwen, 3 epochs & $\Delta$ & Seeds up & 95\% CI \\
\midrule
12 epochs, single model, \Stwo{}$^{\ddagger}$ & \pmsd{0.543}{0.015} & \pmsd{0.462}{0.080} & +0.082$^{\dagger}$ & & \\
All of train, single model, \Stwo{}          & \pmsd{0.492}{0.045} & \pmsd{0.476}{0.043} & +0.017 & 6/10 & \\
All of train, system, \Stwo{}                & 0.575 (0.562) & 0.543 (0.549) & +0.028 & & [$-$0.026, +0.092] \\
\bottomrule
\end{tabular}
\caption{The backbone swap (top) and two follow-ups (bottom) on dev. Single model is the mean\,$\pm$\,s.d.\ over seeds, and $\Delta$ the paired mean difference. LA is logit adjustment. System is the 10-seed ensemble with LA, $\tau$ chosen by nested CV (in parentheses, the mean over 10 CV splits). CIs are bootstraps over dev items, and the 2-annotator sets follow the test regime. $^{\ddagger}$Seeds 0--2. $^{\dagger}$Reported only, since the 12-epoch schedule was adopted on the train slice.}
\label{tab:llm-main}
\end{table*}

\paragraph{Results at ten paired seeds.}
The 3-seed screen passed the bar (+0.081 on \Stwo{}, 3 of 3 seeds up), so we extended the run to ten seeds (Table~\ref{tab:llm-main}).
Per model, Qwen improves \Stwo{} by +0.092 ($t=7.12$) and \Sone{} by +0.095 ($t=8.48$) on all ten seeds, which is more than the whole encoder track gained (0.315 to 0.384).
The system gain holds on the 2-annotator reference sets (0.524 against 0.382).
All the registered predictions about size held, with the \Sone{} gain even exceeding its range, but the prediction about \emph{where} the gain falls turned out wrong (\S\ref{sec:analysis}).
Among published dev results (Figure~\ref{fig:llm}), the system lies above TeleAI's fine-tuned Qwen2.5-7B (0.495) and ChulaNLP's pipeline (0.52) and below TeleAI's multi-call pipeline (0.617), though these come from other protocols.

\FigureLLM

\paragraph{Two single-change follow-ups.}
In the 3-epoch runs the slice F1 rose at every epoch, and the last epoch was always the one selected.
Under a cosine schedule this proves little, because the schedule itself favours the last epoch, so only a longer schedule can test for undertraining.
With 12 epochs (seeds 0--2) the slice selected epochs 9, 11 and 8, and the slice F1 rose by +0.084 on all three seeds, so the rule we had registered adopted the longer schedule.
Training loss fell to about 0.02 without the slice F1 declining, and dev \Stwo{} was \pmsd{0.543}{0.015}, which we report but did not use for the choice.
Training on all 3,448 rows gave +0.017 per model (6 of 10 seeds up), but the system difference has a 95\% interval of [$-$0.026, +0.092], and \Sone{} fell from 0.746 to 0.711.
As with the encoder, the 3-seed screen (+0.039) had overstated the gain.
